
\documentclass{article}
%%%%%%%%%%%%%%%%%%%%%%%%%%%%%%%%%%%%%%%%%%%%%%%%%%%%%%%%%%%%%%%%%%%%%%%%%%%%%%%%%%%%%%%%%%%%%%%%%%%%%%%%%%%%%%%%%%%%%%%%%%%%%%%%%%%%%%%%%%%%%%%%%%%%%%%%%%%%%%%%%%%%%%%%%%%%%%%%%%%%%%%%%%%%%%%%%%%%%%%%%%%%%%%%%%%%%%%%%%%%%%%%%%%%%%%%%%%%%%%%%%%%%%%%%%%%
\usepackage{geometry}
\usepackage{graphicx}
\usepackage{float}

%TCIDATA{OutputFilter=LATEX.DLL}
%TCIDATA{Version=5.50.0.2960}
%TCIDATA{<META NAME="SaveForMode" CONTENT="1">}
%TCIDATA{BibliographyScheme=Manual}
%TCIDATA{Created=Tuesday, April 11, 2023 10:06:45}
%TCIDATA{LastRevised=Thursday, May 04, 2023 08:16:29}
%TCIDATA{<META NAME="GraphicsSave" CONTENT="32">}
%TCIDATA{<META NAME="DocumentShell" CONTENT="Standard LaTeX\Blank - Standard LaTeX Article">}
%TCIDATA{CSTFile=40 LaTeX article.cst}

\newtheorem{theorem}{Theorem}
\newtheorem{acknowledgement}[theorem]{Acknowledgement}
\newtheorem{algorithm}[theorem]{Algorithm}
\newtheorem{axiom}[theorem]{Axiom}
\newtheorem{case}[theorem]{Case}
\newtheorem{claim}[theorem]{Claim}
\newtheorem{conclusion}[theorem]{Conclusion}
\newtheorem{condition}[theorem]{Condition}
\newtheorem{conjecture}[theorem]{Conjecture}
\newtheorem{corollary}[theorem]{Corollary}
\newtheorem{criterion}[theorem]{Criterion}
\newtheorem{definition}[theorem]{Definition}
\newtheorem{example}[theorem]{Example}
\newtheorem{exercise}[theorem]{Exercise}
\newtheorem{lemma}[theorem]{Lemma}
\newtheorem{notation}[theorem]{Notation}
\newtheorem{problem}[theorem]{Problem}
\newtheorem{proposition}[theorem]{Proposition}
\newtheorem{remark}[theorem]{Remark}
\newtheorem{solution}[theorem]{Solution}
\newtheorem{summary}[theorem]{Summary}
\newenvironment{proof}[1][Proof]{\noindent\textbf{#1.} }{\ \rule{0.5em}{0.5em}}
\newcommand{\pwd}{D:/DeepLearning/images_from_K/lesion_normal_comparison/}
\input{tcilatex}
\begin{document}

\title{Spatial Analysis}
\author{}
\maketitle

\section{Methods}

For each sample, we will compare the distribution of each markers in the
normal and lesion regions.

\begin{itemize}
\item Markers for DD: CD3, CD4, CD8a, CD11c, CD15, CD45, CD68, CD69, CD123,
CD279

\item Markers for CF: CD1a, CD3, CD4, CD8a, CD11c, CD15, CD45, CD69, CD123,
CD279 (CD68 did not pass QC)
\end{itemize}

\bigskip

For CD123, segmentation is done at flow=0.5 to get more masks in the bright
area and a threshold of 5000 is set for total intensity to remove masks in
the debris field in the lesion area.

\bigskip

The size of the skin tissue lesion area is determined by morphology using
wheat germ agglutinin (WGA) flourescent stain. ImageJ is used to demarcate
the lesion area, then a straight line ROI is drawn to span from one side of
the lesion to the other side of the lesion, with the endpoints of the line
positioned at the outermost skin boundary. \ 

\bigskip 

Some considerations on region boundary and distance metric:

1. Should we draw the boundary between normal and lesion based on WGA only
or all markers? 

2. Shall each ROI be convex?

For (1), the two choices lead are answering different questions. Drawing
boundary based on WGA only answers a stronger question than drawing the
boundary based on all markers. 

For (2), a convex shape is a shape that a line connecting any two points in
the shape completely lies within the shape. Practically speaking, in this
context it is equivalent to drawing a straight line to separate normal from
lesion. Pros for straight boundary: the span between any cell and the
closest point on the skin lies completely within the region. So the distance
measure is well defined. In contrast, if a line from a cell to its closest
point on the skin surface cuts through both normal and lesion regions, the
interpretation of distance to skin surface is muddled. Pros for curved
boundary: Regions (which is a type of neighborhoods) do not necessarily have
geometric shapes. 

These two questions are also intertwined. If we use WGA only to draw
boundaries, we can only use straight line boundaries because the difference
between normal and lesion is only visible in epidermis in WGA staining. A
straight line can be drawn in epidermis and extend to dermis, but not a
curve.

%TCIMACRO{\TeXButton{clearpage}{\clearpage}}%
%BeginExpansion
\clearpage%
%EndExpansion

\section{Results}

%TCIMACRO{%
%\TeXButton{tab }{\begin{center}
%\begin{table}[h] \centering
%
%\input{\pwd R/tables/input/DD_counts.tex}
%
%\caption{DD}
%\vspace{-5pt} 
%\label{tab DD counts}
%
%\end{table}
%\end{center}}}%
%BeginExpansion
\begin{center}
\begin{table}[h] \centering

\input{\pwd R/tables/input/DD_counts.tex}

\caption{DD}
\vspace{-5pt} 
\label{tab DD counts}

\end{table}
\end{center}%
%EndExpansion

%TCIMACRO{%
%\TeXButton{tab }{\begin{center}
%\begin{table}[h] \centering
%
%\input{\pwd R/tables/input/CF_counts.tex}
%
%\caption{CF}
%\vspace{-5pt} 
%\label{tab DD counts}
%
%\end{table}
%\end{center}}}%
%BeginExpansion
\begin{center}
\begin{table}[h] \centering

\input{\pwd R/tables/input/CF_counts.tex}

\caption{CF}
\vspace{-5pt} 
\label{tab DD counts}

\end{table}
\end{center}%
%EndExpansion

%TCIMACRO{%
%\TeXButton{tab }{\begin{center}
%\begin{table}[h] \centering
%
%\input{\pwd R/tables/input/peaks_distance_lesion_size.tex}
%
%\caption{Comparing the space between CD3 and CD15 peaks with lesion size. Distance is measured in um, size is measured in square mm.}
%\vspace{-5pt} 
%\label{tab DD counts}
%
%\end{table}
%\end{center}}}%
%BeginExpansion
\begin{center}
\begin{table}[h] \centering

\input{\pwd R/tables/input/peaks_distance_lesion_size.tex}

\caption{Comparing the space between CD3 and CD15 peaks with lesion size. Distance is measured in um, size is measured in square mm.}
\vspace{-5pt} 
\label{tab DD counts}

\end{table}
\end{center}%
%EndExpansion

%TCIMACRO{\TeXButton{clearpage}{\clearpage}}%
%BeginExpansion
\clearpage%
%EndExpansion

%TCIMACRO{%
%\TeXButton{fig}{\begin{center}
%\begin{figure}[H] \centering
%\includegraphics[width=\textwidth]{\pwd R/figures/compare_DD_CF_dist2skin_count.png}
%
%\caption{Distance to skin, counts.}
%\vspace{-5pt}
%\label{}
%\end{figure}
%\end{center}}}%
%BeginExpansion
\begin{center}
\begin{figure}[H] \centering
\includegraphics[width=\textwidth]{\pwd R/figures/compare_DD_CF_dist2skin_count.png}

\caption{Distance to skin, counts.}
\vspace{-5pt}
\label{}
\end{figure}
\end{center}%
%EndExpansion

%TCIMACRO{%
%\TeXButton{fig}{\begin{center}
%\begin{figure}[H] \centering
%\includegraphics[width=\textwidth]{\pwd R/figures/compare_DD_CF_dist2DEJ_count.png}
%
%\caption{Distance to DEJ, counts.}
%\vspace{-5pt}
%\label{}
%\end{figure}
%\end{center}}}%
%BeginExpansion
\begin{center}
\begin{figure}[H] \centering
\includegraphics[width=\textwidth]{\pwd R/figures/compare_DD_CF_dist2DEJ_count.png}

\caption{Distance to DEJ, counts.}
\vspace{-5pt}
\label{}
\end{figure}
\end{center}%
%EndExpansion

%TCIMACRO{%
%\TeXButton{fig}{\begin{center}
%\begin{figure}[H] \centering
%\includegraphics[width=\textwidth]{\pwd R/figures/compare_DD_CF_dist2skin_density.png}
%
%\caption{Distance to skin, density.}
%\vspace{-5pt}
%\label{}
%\end{figure}
%\end{center}}}%
%BeginExpansion
\begin{center}
\begin{figure}[H] \centering
\includegraphics[width=\textwidth]{\pwd R/figures/compare_DD_CF_dist2skin_density.png}

\caption{Distance to skin, density.}
\vspace{-5pt}
\label{}
\end{figure}
\end{center}%
%EndExpansion

%TCIMACRO{%
%\TeXButton{fig}{\begin{center}
%\begin{figure}[H] \centering
%\includegraphics[width=\textwidth]{\pwd R/figures/compare_DD_CF_dist2DEJ_density.png}
%
%\caption{Distance to DEJ, density.}
%\vspace{-5pt}
%\label{}
%\end{figure}
%\end{center}}}%
%BeginExpansion
\begin{center}
\begin{figure}[H] \centering
\includegraphics[width=\textwidth]{\pwd R/figures/compare_DD_CF_dist2DEJ_density.png}

\caption{Distance to DEJ, density.}
\vspace{-5pt}
\label{}
\end{figure}
\end{center}%
%EndExpansion

%TCIMACRO{%
%\TeXButton{fig}{\begin{center}
%\begin{figure}[H] \centering
%\includegraphics[width=.45\textwidth]{\pwd R/figures/DD_dist2DEJ_compare_markers.png}
%\includegraphics[width=.45\textwidth]{\pwd R/figures/DD_dist2skin_compare_markers.png}
%
%\caption{Distance to DEJ (left) and skin surface (right).}
%\vspace{-5pt}
%\label{}
%\end{figure}
%\end{center}}}%
%BeginExpansion
\begin{center}
\begin{figure}[H] \centering
\includegraphics[width=.45\textwidth]{\pwd R/figures/DD_dist2DEJ_compare_markers.png}
\includegraphics[width=.45\textwidth]{\pwd R/figures/DD_dist2skin_compare_markers.png}

\caption{Distance to DEJ (left) and skin surface (right).}
\vspace{-5pt}
\label{}
\end{figure}
\end{center}%
%EndExpansion

%TCIMACRO{%
%\TeXButton{fig}{\begin{center}
%\begin{figure}[H] \centering
%\includegraphics[width=.45\textwidth]{\pwd R/figures/CF_dist2DEJ_compare_markers.png}
%\includegraphics[width=.45\textwidth]{\pwd R/figures/CF_dist2skin_compare_markers.png}
%
%\caption{Distance to DEJ (left) and skin surface (right).}
%\vspace{-5pt}
%\label{}
%\end{figure}
%\end{center}}}%
%BeginExpansion
\begin{center}
\begin{figure}[H] \centering
\includegraphics[width=.45\textwidth]{\pwd R/figures/CF_dist2DEJ_compare_markers.png}
\includegraphics[width=.45\textwidth]{\pwd R/figures/CF_dist2skin_compare_markers.png}

\caption{Distance to DEJ (left) and skin surface (right).}
\vspace{-5pt}
\label{}
\end{figure}
\end{center}%
%EndExpansion

\bigskip 

\section{Appendix}

\end{document}
